% Options for packages loaded elsewhere
% Options for packages loaded elsewhere
\PassOptionsToPackage{unicode}{hyperref}
\PassOptionsToPackage{hyphens}{url}
\PassOptionsToPackage{dvipsnames,svgnames,x11names}{xcolor}
%
\documentclass[
  12pt,
  letterpaper,
]{article}
\usepackage{xcolor}
\usepackage[margin=1in]{geometry}
\usepackage{amsmath,amssymb}
\setcounter{secnumdepth}{5}
\usepackage{iftex}
\ifPDFTeX
  \usepackage[T1]{fontenc}
  \usepackage[utf8]{inputenc}
  \usepackage{textcomp} % provide euro and other symbols
\else % if luatex or xetex
  \usepackage{unicode-math} % this also loads fontspec
  \defaultfontfeatures{Scale=MatchLowercase}
  \defaultfontfeatures[\rmfamily]{Ligatures=TeX,Scale=1}
\fi
\usepackage{lmodern}
\ifPDFTeX\else
  % xetex/luatex font selection
\fi
% Use upquote if available, for straight quotes in verbatim environments
\IfFileExists{upquote.sty}{\usepackage{upquote}}{}
\IfFileExists{microtype.sty}{% use microtype if available
  \usepackage[]{microtype}
  \UseMicrotypeSet[protrusion]{basicmath} % disable protrusion for tt fonts
}{}
\makeatletter
\@ifundefined{KOMAClassName}{% if non-KOMA class
  \IfFileExists{parskip.sty}{%
    \usepackage{parskip}
  }{% else
    \setlength{\parindent}{0pt}
    \setlength{\parskip}{6pt plus 2pt minus 1pt}}
}{% if KOMA class
  \KOMAoptions{parskip=half}}
\makeatother
% Make \paragraph and \subparagraph free-standing
\makeatletter
\ifx\paragraph\undefined\else
  \let\oldparagraph\paragraph
  \renewcommand{\paragraph}{
    \@ifstar
      \xxxParagraphStar
      \xxxParagraphNoStar
  }
  \newcommand{\xxxParagraphStar}[1]{\oldparagraph*{#1}\mbox{}}
  \newcommand{\xxxParagraphNoStar}[1]{\oldparagraph{#1}\mbox{}}
\fi
\ifx\subparagraph\undefined\else
  \let\oldsubparagraph\subparagraph
  \renewcommand{\subparagraph}{
    \@ifstar
      \xxxSubParagraphStar
      \xxxSubParagraphNoStar
  }
  \newcommand{\xxxSubParagraphStar}[1]{\oldsubparagraph*{#1}\mbox{}}
  \newcommand{\xxxSubParagraphNoStar}[1]{\oldsubparagraph{#1}\mbox{}}
\fi
\makeatother


\usepackage{longtable,booktabs,array}
\newcounter{none} % for unnumbered tables
\usepackage{calc} % for calculating minipage widths
% Correct order of tables after \paragraph or \subparagraph
\usepackage{etoolbox}
\makeatletter
\patchcmd\longtable{\par}{\if@noskipsec\mbox{}\fi\par}{}{}
\makeatother
% Allow footnotes in longtable head/foot
\IfFileExists{footnotehyper.sty}{\usepackage{footnotehyper}}{\usepackage{footnote}}
\makesavenoteenv{longtable}
\usepackage{graphicx}
\makeatletter
\newsavebox\pandoc@box
\newcommand*\pandocbounded[1]{% scales image to fit in text height/width
  \sbox\pandoc@box{#1}%
  \Gscale@div\@tempa{\textheight}{\dimexpr\ht\pandoc@box+\dp\pandoc@box\relax}%
  \Gscale@div\@tempb{\linewidth}{\wd\pandoc@box}%
  \ifdim\@tempb\p@<\@tempa\p@\let\@tempa\@tempb\fi% select the smaller of both
  \ifdim\@tempa\p@<\p@\scalebox{\@tempa}{\usebox\pandoc@box}%
  \else\usebox{\pandoc@box}%
  \fi%
}
% Set default figure placement to htbp
\def\fps@figure{htbp}
\makeatother


% definitions for citeproc citations
\NewDocumentCommand\citeproctext{}{}
\NewDocumentCommand\citeproc{mm}{%
  \begingroup\def\citeproctext{#2}\cite{#1}\endgroup}
\makeatletter
 % allow citations to break across lines
 \let\@cite@ofmt\@firstofone
 % avoid brackets around text for \cite:
 \def\@biblabel#1{}
 \def\@cite#1#2{{#1\if@tempswa , #2\fi}}
\makeatother
\newlength{\cslhangindent}
\setlength{\cslhangindent}{1.5em}
\newlength{\csllabelwidth}
\setlength{\csllabelwidth}{3em}
\newenvironment{CSLReferences}[2] % #1 hanging-indent, #2 entry-spacing
 {\begin{list}{}{%
  \setlength{\itemindent}{0pt}
  \setlength{\leftmargin}{0pt}
  \setlength{\parsep}{0pt}
  % turn on hanging indent if param 1 is 1
  \ifodd #1
   \setlength{\leftmargin}{\cslhangindent}
   \setlength{\itemindent}{-1\cslhangindent}
  \fi
  % set entry spacing
  \setlength{\itemsep}{#2\baselineskip}}}
 {\end{list}}
\usepackage{calc}
\newcommand{\CSLBlock}[1]{\hfill\break\parbox[t]{\linewidth}{\strut\ignorespaces#1\strut}}
\newcommand{\CSLLeftMargin}[1]{\parbox[t]{\csllabelwidth}{\strut#1\strut}}
\newcommand{\CSLRightInline}[1]{\parbox[t]{\linewidth - \csllabelwidth}{\strut#1\strut}}
\newcommand{\CSLIndent}[1]{\hspace{\cslhangindent}#1}



\setlength{\emergencystretch}{3em} % prevent overfull lines

\providecommand{\tightlist}{%
  \setlength{\itemsep}{0pt}\setlength{\parskip}{0pt}}



 


\usepackage{booktabs}
\usepackage{longtable}
\usepackage{array}
\usepackage{multirow}
\usepackage{wrapfig}
\usepackage{float}
\usepackage{colortbl}
\usepackage{pdflscape}
\usepackage{tabu}
\usepackage{threeparttable}
\usepackage{threeparttablex}
\usepackage[normalem]{ulem}
\usepackage{makecell}
\usepackage{xcolor}
\usepackage{caption}
\usepackage{anyfontsize}
\usepackage[bottom]{footmisc}
\usepackage{setspace}
\usepackage{parskip}
\setlength{\parindent}{20pt}
\usepackage{sectsty}
\sectionfont{\fontsize{12}{15}\selectfont}
\subsectionfont{\fontsize{12}{15}\selectfont}
\usepackage{fancyhdr}
\fancyhead{}
%\fancyhead[L]{\small\textit{Work in progress. Please do not circulate.}}
\pagestyle{fancy}
\renewcommand{\headrulewidth}{0pt}
\usepackage{tocloft}
\renewcommand{\cftsecdotsep}{\cftdotsep}
\renewcommand{\thefigure}{S\arabic{figure}}
\renewcommand{\thetable}{S\arabic{table}}
\renewcommand{\figurename}{Fig.}
\renewcommand{\tablename}{Table}
\makeatletter
\@ifpackageloaded{tcolorbox}{}{\usepackage[skins,breakable]{tcolorbox}}
\@ifpackageloaded{fontawesome5}{}{\usepackage{fontawesome5}}
\definecolor{quarto-callout-color}{HTML}{909090}
\definecolor{quarto-callout-note-color}{HTML}{0758E5}
\definecolor{quarto-callout-important-color}{HTML}{CC1914}
\definecolor{quarto-callout-warning-color}{HTML}{EB9113}
\definecolor{quarto-callout-tip-color}{HTML}{00A047}
\definecolor{quarto-callout-caution-color}{HTML}{FC5300}
\definecolor{quarto-callout-color-frame}{HTML}{acacac}
\definecolor{quarto-callout-note-color-frame}{HTML}{4582ec}
\definecolor{quarto-callout-important-color-frame}{HTML}{d9534f}
\definecolor{quarto-callout-warning-color-frame}{HTML}{f0ad4e}
\definecolor{quarto-callout-tip-color-frame}{HTML}{02b875}
\definecolor{quarto-callout-caution-color-frame}{HTML}{fd7e14}
\makeatother
\makeatletter
\@ifpackageloaded{caption}{}{\usepackage{caption}}
\AtBeginDocument{%
\ifdefined\contentsname
  \renewcommand*\contentsname{Table of contents}
\else
  \newcommand\contentsname{Table of contents}
\fi
\ifdefined\listfigurename
  \renewcommand*\listfigurename{List of Figures}
\else
  \newcommand\listfigurename{List of Figures}
\fi
\ifdefined\listtablename
  \renewcommand*\listtablename{List of Tables}
\else
  \newcommand\listtablename{List of Tables}
\fi
\ifdefined\figurename
  \renewcommand*\figurename{Figure}
\else
  \newcommand\figurename{Figure}
\fi
\ifdefined\tablename
  \renewcommand*\tablename{Table}
\else
  \newcommand\tablename{Table}
\fi
}
\@ifpackageloaded{float}{}{\usepackage{float}}
\floatstyle{ruled}
\@ifundefined{c@chapter}{\newfloat{codelisting}{h}{lop}}{\newfloat{codelisting}{h}{lop}[chapter]}
\floatname{codelisting}{Listing}
\newcommand*\listoflistings{\listof{codelisting}{List of Listings}}
\makeatother
\makeatletter
\makeatother
\makeatletter
\@ifpackageloaded{caption}{}{\usepackage{caption}}
\@ifpackageloaded{subcaption}{}{\usepackage{subcaption}}
\makeatother
\usepackage{bookmark}
\IfFileExists{xurl.sty}{\usepackage{xurl}}{} % add URL line breaks if available
\urlstyle{same}
\hypersetup{
  pdftitle={Supplementary materials for Greening while Governing the Economy: Do Green Central Banks Matter?},
  colorlinks=true,
  linkcolor={blue},
  filecolor={Maroon},
  citecolor={Blue},
  urlcolor={blue},
  pdfcreator={LaTeX via pandoc}}


\title{Supplementary materials for\\
Greening while Governing the Economy: Do Green Central Banks Matter?}
\author{}
\date{}
\begin{document}
\maketitle

\renewcommand*\contentsname{Table of Contents}
{
\setcounter{tocdepth}{3}
\tableofcontents
}

\newpage{}

\section{Materials}\label{materials}

\subsection{Dependent Variable: Greenhouse Gas
Emissions}\label{dependent-variable-greenhouse-gas-emissions}

My dependent variable consists of greenhouse gas (GHG) emissions. The
data have been collected from the Emissions Database for Global
Atmospheric Research (EDGAR) developed by the European Commission's
Joint Research Centre (\emph{1}). The dataset is publicly available
(\url{https://edgar.jrc.ec.europa.eu/dataset_ghg2025}) and has been
widely used (e.g., \emph{2}, \emph{3}). EDGAR has separate estimates for
countries' total carbon dioxide (\(\text{CO}_2\)), methane
(\(\text{CH}_4\)), nitrous oxide (\(\text{N}_2\text{O}\)), and
fluorinated gases, as well as total GHG emissions measured in millions
of \(\text{CO}_2\) equivalent, which is the one I use in this article.
From the original spreadsheets downloaded from EDGAR, I extracted the
total GHG emissions data for the years 2001 to 2021.

EDGAR details its methodology to calculate emissions here:
\url{https://edgar.jrc.ec.europa.eu/methodology}. Specifically, EDGAR
quantifies global anthropogenic emissions by applying a consistent
estimation framework across all countries and economic sectors by year.
This approach ensures a highly comparable data baseline across
heterogeneous regulatory regimes and distinct gas categories. Only two
countries needed to be excluded from the assessment due to missing data
on the dependent variable: Marshall Islands and Northern Mariana
Islands.

I use GHG emissions as a dependent variable based on three
specifications. First, given the skewed distribution of GHG emissions
across countries and years, I transformed this variable to a logarithmic
format. Additionally, I normalized GHG emissions by the log of
countries' total population to obtain a measure of emissions per capita.
I also normalized GHG emissions by the log of gross domestic product
(GDP) per capita (PC) to account for economic and consumption
differences across countries. The normalization process involves taking
the logarithm of the ratio of GHG emissions to the population or GDP PC.
An example of emissions per capita is expressed mathematically as:

\begin{equation}\protect\phantomsection\label{eq-normalization}{
\log(ghgpc_{i,t}) = \log(ghg_{i,t}) - \log(pop_{i,t})
}\end{equation}

where \(\log(ghgpc_{i,t})\) is the log-transformed emissions per capita
variable for country \(i\) in year \(t\). \(\log(ghg_{i,t})\) is the log
of GHG emissions. \(\log(pop_{i,t})\) is the normalization variable in
its log-transformed version. In the Equation~\ref{eq-normalization}
example, I use population, but the same process applies to GDP per
capita. I explain from which repositories the other variables have been
collected and how I have built a merged dataset below.

\subsection{Independent and Control
Variables}\label{independent-and-control-variables}

\subsubsection{NGFS}\label{ngfs}

My main independent variable is membership in the Network for Greening
the Financial System (NGFS), which serves as a treatment indicator of
financial regulators' commitment to incorporate climate risks into
financial supervision. The NGFS maintains their membership directory
updated and publicly available here:
\url{https://www.ngfs.net/en/about-us/membership}. However, the NGFS
does not specify the year in which financial supervisors have become
members or observers. In March 2024, NGFS officers provided a
spreadsheet containing the years in which each member and observer
joined the network until then. Since this is not available through the
website above, the e-mail exchange and the original spreadsheet shared
by the NGFS are available upon request for reviewers and editors. I
excluded observers and restricted the NGFS data to members only between
2001 and 2021.

Subsequently, I identified each country where financial supervisors are
located and recorded the year a nation joined the NGFS. This way,
countries are not double counted. Abu Dhabi, Dubai, and the European
Union (EU) were also removed from the data. Abu Dhabi and Dubai count as
``United Arab Emirates'', whereas several EU nations appear
independently in the dataset. I then created a dummy variable,
\texttt{ngfs}, indicating NGFS membership and assigned a value of 1 to
countries in the year in which they became NGFS members and 0 to
non-members. Thus, non-members include nations \emph{before} they joined
the NGFS and countries that never became members within the period
analyzed. Consider, for example, how France has been coded in
Table~\ref{tbl-france-ngfs}, as the Banque de France helped found the
NGFS in 2017 (\emph{4}).

\newpage{}

\begin{longtable}[t]{ccc}

\caption{\label{tbl-france-ngfs}Example of NGFS Membership Coding:
France (2015--2021)}

\tabularnewline

\toprule
Country & Year & NGFS Member\\
\midrule
France & 2001 & 0\\
France & 2002 & 0\\
France & 2003 & 0\\
France & 2004 & 0\\
France & 2005 & 0\\
\addlinespace
France & 2006 & 0\\
France & 2007 & 0\\
France & 2008 & 0\\
France & 2009 & 0\\
France & 2010 & 0\\
\addlinespace
France & 2011 & 0\\
France & 2012 & 0\\
France & 2013 & 0\\
France & 2014 & 0\\
France & 2015 & 0\\
\addlinespace
France & 2016 & 0\\
France & 2017 & 1\\
France & 2018 & 1\\
France & 2019 & 1\\
France & 2020 & 1\\
\addlinespace
France & 2021 & 1\\
\bottomrule

\end{longtable}

As a result of the steps above, I built an initial
\(country \times year\) dataset, which I merged with other repositories
from which other variables were obtained. In sum, the cross-sectional
unit \(i\) is a country, which appears in each year \(t\) over a period
of time. I detail other methodological decisions that resulted in the
tidy, merged, final dataset used for the analysis in
Section~\ref{sec-methods}.

\subsubsection{Central Bank Independence
Index}\label{central-bank-independence-index}

A second independent variable of interest is central bank independence
(CBI) (\emph{5}, \emph{6}), which is publicly available
(\url{https://sites.google.com/site/carogarriga/cbi-data-1}). CBI is an
index ranging from 0 to 1, where 0 represents no independence and 1
denotes maximum independence (\emph{5}). The index considers four formal
attributes of central banks: (i) personnel independence of a central
bank's executive leadership; (ii) having exclusive power or not over
price stability mandate in terms of policy objectives; (iii)
independence in policy formulation to devise rules, norms, and
regulations; and (iv) lending limitations when furnishing money to the
own government (\emph{5}). These metrics are compiled into two
operational variables: an unweighted index that computes the simple
average of all four components; and a weighted index that mathematically
penalizes loose sovereign lending rules by weighting restriction
variables more heavily.

I use Garriga's (\emph{5}) weighted CBI index for two reasons. First, it
``represents a 46 percent extension of the coverage of the most
comprehensive dataset on de jure CBI available to date (Garriga 2016).
On average, this dataset includes fifty four more countries per year
than Bodea and Hicks (2015b, updated to 2020), thirty seven more than
Romelli (2022),and twenty two more than Garriga (2016)'' (\emph{5}). And
second, it offers more conservative estimates than other CBI datasets:
``Due to systematically missing data on countries for whom information
is harder to obtain, previous research may have overestimated the level
of CBI in the world, which may have biased inferences based on smaller
samples'' (\emph{5}). Therefore, this dataset provides more data on
developing nations over a longer period of time than other repositories,
and it helps address potential issues with overestimating CBI. As a
robustness check, I also tested the unweighted index and found similar
results (see Section~\ref{sec-robustness}).

\subsubsection{World Bank}\label{world-bank}

Following Chesler et al. (\emph{7}), a set of control variables includes
data from the World Bank, which have been collected using the
\texttt{wbstats} package in \texttt{R} (\emph{8}). To normalize the
distribution of the controls either as a logarithm or square root, they
were transformed based on the closest to a normal distribution displayed
using the \texttt{gladder} function of the \texttt{describedata} package
also in \texttt{R} (\emph{9}). Essentially, \texttt{gladder} plots a
series of potential variable transformations while showing a normal
distribution line on each transformation. The indicators are the
following.

\textbf{Log of GDP per capita} accounts for a country's standard of
living and its population purchasing power, which has long been
positively correlated with higher emissions (\emph{7}). \textbf{Log of
GDP per capita growth}, in turn, measures economic growth and whether
the income level per person has increased or decreased, and it also
reflects an increase in emissions (\emph{7}). \textbf{Log of population
growth} follows demographic trends regarding population size, which
yields inconclusive results for emissions when controlling for other
factors like GDP per capita.\footnote{This has been a long debate that
  traces back to Ostrom's (\emph{10}) work on challenging the assumption
  that population growth inherently leads to the depletion of natural
  resources.} \textbf{Log of total natural resources rented} is likely
positively associated with higher emissions (\emph{7}), and it consists
of the sum of, e.g., oil, forests, and mineral rents, among other
assets, rented as a percentage of a nation's GDP. \textbf{Log of trade
dependency} is measured as the ratio of a country's total trade to GDP,
and the circulation of products tends to increase emissions (\emph{7}).
\textbf{Square root of agricultural land} refers to arable areas used
either for temporary or permanent crops as a percentage of total land,
which may also produce higher emission levels (\emph{7}). And,
conversely, the \textbf{square root of forest land} refers to areas
natural or planted trees as a percentage of total land, which might help
offset emissions increased by land use and land cover change.

\subsubsection{Varieties of Democracy}\label{varieties-of-democracy}

Also drawing from Chesler et al. (\emph{7}), as well as Kakenmaster
(\emph{2}), political indicators were collected from the Varieties of
Democracy (VDM) project using the \texttt{vdem} package in \texttt{R}
(\emph{11}, \emph{12}). The VDEM variables are the following.
\textbf{Political corruption index} (\texttt{v2x\_corr}) ranges from 0
to 1, corresponding to the less corrupt to the more corrupt nations.
Corruption levels shape the relationship between political regimes and
GHG emissions (\emph{2}, \emph{13}), with growing levels of corruption
being associated with higher emissions. \textbf{Deliberative democracy
index} (\texttt{v2x\_delibdem}) also ranges from 0 to 1, with higher
values denoting more effective deliberative democratic processes and a
focus on common goods. As deliberations are key in climate and
environmental governance (\emph{10}, \emph{14}), it is expected a
negative relationship between deliberative democracy values and GHG
emissions. Finally, following Vaccaro (\emph{15}), I select a variable
of \textbf{state capacity} (\texttt{v2clrspct}) that captures the extent
to which state agents can abide to the law while pursuing their mandates
while arbitrariness and biased public administration are limited. I
employed min-max re-scaling for state capacity to range from 0 to
1,\footnote{The min-max normalization is mathematically expressed as \[
  \text{State Capacity}_{i,t}^{\text{scaled}} = \frac{\text{State Capacity}_{i,t} - \min(\text{State Capacity})}{\max(\text{State Capacity}) - \min(\text{State Capacity})}
  \] where \(\text{State Capacity}_{i,t}\) represents the raw state
  capacity score for country \(i\) in year \(t\), and
  \(\min(\text{State Capacity})\) and \(\max(\text{State Capacity})\)
  are the sample minimum and maximum values, respectively.} such as the
other VDEM indicators. The higher the score, the higher the levels of
capacity, which, in turn, are presumed to be associated with the ability
to integrate financial and climate policies.

\section{Methodology}\label{sec-methods}

Methodological decisions often entail considering trade-offs,
particularly with respect to dropping observations and selecting a
timeframe for analysis. Here, the example displayed in
Table~\ref{tbl-france-ngfs} raises two main issues. First, by creating a
dataset spanning 2001 and 2024, the pre-treatment period---before
2017---has a disproportionate number of years for which a country being
an NGFS member equals 0. This issue could overestimate the magnitude of
the treatment effect.

Second, several control variables obtained from the World Bank had
missing data for the years between 2021 and 2024, when I received the
NGFS membership data. Because this is one of the first studies relying
on NGFS participation as a treatment indicator using real-world data
with observed values, I have prioritized to be as conservative as
possible with missing data. Hence, I use listwise deletion. As Pepinsky
(\emph{16}) notes, ``Only knowing the true values from the fully
observed data allows us to compare the performance of {[}listwise
deletion and multiple imputation{]}. But this exercise clearly reveals
that multiple imputation can change substantive conclusions in
real-world applications, and in this case, those conclusions will
frequently be more misleading than those that come from listwise
deletion.'' I acknowledge that this approach might present issues also
discussed by Pepinsky (\emph{16}). However, since I use indicators based
on real-world observed values and I aim to estimate conservative
coefficients employing causal inference techniques, my main goal has
been to work using a balanced panel dataset.

To deal with the first issue above, I further restricted the initial
dataset to the period between 2015---almost immediately before the
establishment of the NGFS in 2017---and 2021. And, to address the second
issue, I ensure that each country appears in every single year with no
missing data on the dependent and independent variables. These steps
result in a balanced panel dataset that is suitable for both standard
TWFE modeling, as well as subsequent and more robust
difference-in-differences analysis. The DID with TWFE models are based
on a \(country \times year\) sample of 133 nations that appear over 7
years, resulting in \(n = 931\).

\subsection{Descriptive Statistics}\label{descriptive-statistics}

For an overview of the data, I break down the descriptive statistics by
NGFS membership. Table~\ref{tbl-desc-stats} highlights the differences
between groups. This includes nations before they received the
treatment, after the treatment, and countries that never joined the
NGFS.

\begin{table}[H]

\caption{\label{tbl-desc-stats}Descriptive statistics by NGFS membership
treatment (N = 931)}

\centering{

\fontsize{12.0pt}{14.0pt}\selectfont
\begin{tabular*}{\linewidth}{@{\extracolsep{\fill}}>{\raggedright\arraybackslash}p{\dimexpr 0.45\linewidth -2\tabcolsep-1.5\arrayrulewidth}>{\centering\arraybackslash}p{\dimexpr 0.20\linewidth -2\tabcolsep-1.5\arrayrulewidth}>{\centering\arraybackslash}p{\dimexpr 0.20\linewidth -2\tabcolsep-1.5\arrayrulewidth}>{\centering\arraybackslash}p{\dimexpr 0.15\linewidth -2\tabcolsep-1.5\arrayrulewidth}}
\toprule
\textbf{Variable} & \multicolumn{1}{>{\centering\arraybackslash}m{\dimexpr 0.20\linewidth -2\tabcolsep-1.5\arrayrulewidth}}{{\shortstack[c]{\parbox{\linewidth}{\centering \textbf{Member}  \\ N = 193}}}} & \multicolumn{1}{>{\centering\arraybackslash}m{\dimexpr 0.20\linewidth -2\tabcolsep-1.5\arrayrulewidth}}{{\shortstack[c]{\parbox{\linewidth}{\centering \textbf{Non-Member}  \\ N = 738}}}} & \multicolumn{1}{>{\centering\arraybackslash}m{\dimexpr 0.15\linewidth -2\tabcolsep-1.5\arrayrulewidth}}{{\shortstack[c]{\parbox{\linewidth}{\centering \textbf{Differences}  \\ p-value}}}} \\ 
\midrule\addlinespace[2.5pt]
Greenhouse gas emissions & 11.83 (1.82) & 10.83 (1.74) & {\bfseries <0.001} \\ 
Greenhouse gas emissions per GDP per capita & 1.97 (2.16) & 2.40 (1.85) & {\bfseries 0.004} \\ 
Greenhouse gas emissions per capita & -4.87 (0.59) & -5.44 (0.95) & {\bfseries <0.001} \\ 
Trade (\% of GDP) & 4.41 (0.59) & 4.25 (0.49) & {\bfseries 0.005} \\ 
GDP per capita (constant 2015 US\$) & 9.86 (1.07) & 8.43 (1.37) & {\bfseries <0.001} \\ 
GDP per capita growth (annual \%) & 3.93 (0.11) & 3.94 (0.10) & 0.5 \\ 
Natural resources & -0.86 (2.46) & 0.60 (2.06) & {\bfseries <0.001} \\ 
Population growth (annual \%) & 2.52 (0.07) & 2.58 (0.14) & {\bfseries <0.001} \\ 
Agricultural land (\% of land area) & 5.83 (2.00) & 6.19 (1.82) & 0.082 \\ 
Forest area (\% of land area) & 5.54 (1.83) & 4.93 (2.33) & {\bfseries 0.003} \\ 
Deliberative democracy index (0-1) & 0.60 (0.24) & 0.40 (0.23) & {\bfseries <0.001} \\ 
Central Bank Independence Index (Weighted) & 0.66 (0.20) & 0.62 (0.19) & {\bfseries <0.001} \\ 
Region &  &  &  \\ 
    East and South Asia and Pacific & 40 (21\%) & 107 (14\%) &  \\ 
    Europe and Central Asia & 95 (49\%) & 213 (29\%) &  \\ 
    Latin America and Caribbean & 23 (12\%) & 103 (14\%) &  \\ 
    Middle East and Africa & 29 (15\%) & 307 (42\%) &  \\ 
    North America & 6 (3.1\%) & 8 (1.1\%) &  \\ 
\bottomrule
\end{tabular*}
\begin{minipage}{\linewidth}
\vspace{.05em}
The values correspond to the mean for continuous
variables and the proportion for categorical variables. The standard
deviation for the continuous variables is shown in parentheses.\\
\end{minipage}

}

\end{table}%

Table~\ref{tbl-desc-stats} reveals some significant differences worth
noting. NGFS members (\(n = 193\)) emit substantially more GHG in
absolute terms than non-members (\(n = 738\)) (mean log GHG = 11.83
vs.~10.83, \(p < 0.001\)). This gap reflects the NGFS structure: Early
adopters are predominantly industrialized economies whose central banks
have the capacity and political motivation to engage with international
financial governance (\emph{4}). Compared to non-members, NGFS members
are notably wealthier (log GDP per capita: 9.86 vs.~8.43,
\(p < 0.001\)), more trade-exposed (4.41 vs.~4.25, \(p = 0.005\)), and
score higher on deliberative democracy (0.60 vs.~0.40, \(p < 0.001\)).
At the same time, emissions per GDP per capita are lower among NGFS
members as opposed to non-members (1.97 vs.~2.40, \(p = 0.004\)). These
systematic differences underscore that standard before-after or
cross-sectional models may produce unreliable estimates of the NGFS
effect on emissions, motivating the within-country identification
strategy I employ here.

Indeed, the regional distribution of NGFS membership reinforces this
concern. Europe and Central Asia account for nearly half of all
member-country observations (49\%) but less than a third of non-member
observations (29\%). By contrast, Middle East and Africa, which
represent 42\% of non-member observations, represent only 15\% of NGFS
member observations. Yet, GDP per capita growth does not differ
significantly between members and non-members (3.93 vs.~3.94,
\(p = 0.5\)). This similarity in growth trajectories matters: It
suggests that the emissions reduction are not explained by members
simply growing more slowly than non-members. Instead, it is consistent
with pathways to decarbonization rather than economic deceleration.
Relatedly, central bank independence is slightly but significantly
higher among members (0.66 vs.~0.62, \(p < 0.001\)). This pattern
suggests that central banks need some level of statutory insulation from
short-term political interference before they can extend their monetary
mandate to design policies targeted at handling environmental
externalities.

The differences in Table~\ref{tbl-desc-stats} have both statistical and
practical implications. Substantively, NGFS membership has thus far been
concentrated among nations whose central banks already operate in
well-developed regulatory environments. This feature may amplify the
conditional effects related to the independence of financial
supervisors. Empirically, the variation in CBI among NGFS members and
non-members is modest in magnitude but important to further test whether
and how CBI works as a moderator or mediation variable in the models
discussed in Section~\ref{sec-models}. Finally, the sample differences
raise questions about whether the findings generalize beyond the current
NGFS membership composition, which I address by conducting robustness
checks (see Section~\ref{sec-robustness}).

\begin{tcolorbox}[enhanced jigsaw, arc=.35mm, bottomrule=.15mm, bottomtitle=1mm, breakable, colback=white, colbacktitle=quarto-callout-note-color!10!white, colframe=quarto-callout-note-color-frame, coltitle=black, left=2mm, leftrule=.75mm, opacityback=0, opacitybacktitle=0.6, rightrule=.15mm, title=\textcolor{quarto-callout-note-color}{\faInfo}\hspace{0.5em}{Note}, titlerule=0mm, toprule=.15mm, toptitle=1mm]

I need to double check how best to conduct the robustness check above on
sample composition.

\end{tcolorbox}

\subsection{Regression Results}\label{sec-models}

\section{Robustness Checks}\label{sec-robustness}

\section{Replication Notes}\label{replication-notes}

The documents and code to replicate the results are available upon
request, and they will be published as part of a GitHub public
repository.

\newpage{}

\section*{References}\label{references}
\addcontentsline{toc}{section}{References}

\protect\phantomsection\label{refs}
\begin{CSLReferences}{0}{1}
\bibitem[\citeproctext]{ref-crippa2025}
\CSLLeftMargin{1. }%
\CSLRightInline{M. Crippa, D. Guizzardi, F. Pagani, M. Banja, M.
Muntean, {``{GHG Emissions} of {All World Countries} - 2025 {Report}''}
(JRC143227, Publications Office of the European Union, Luxembourg,
2025); \url{https://doi.org/10.2760/9816914}.}

\bibitem[\citeproctext]{ref-kakenmaster2025}
\CSLLeftMargin{2. }%
\CSLRightInline{W. Kakenmaster,
\href{https://doi.org/10.1017/S1537592724000793}{The {Fossil-Fueled
Roots} of {Climate Inaction} in {Authoritarian Regimes}}.
\emph{Perspectives on Politics} \textbf{23}, 415--433 (2025).}

\bibitem[\citeproctext]{ref-stechemesser2024}
\CSLLeftMargin{3. }%
\CSLRightInline{A. Stechemesser, N. Koch, E. Mark, E. Dilger, P. Klösel,
L. Menicacci, D. Nachtigall, F. Pretis, N. Ritter, M. Schwarz, H.
Vossen, A. Wenzel,
\href{https://doi.org/10.1126/science.adl6547}{Climate {Policies That
Achieved Major Emission Reductions}: {Global Evidence} from {Two
Decades}}. \emph{Science} \textbf{385}, 884--892 (2024).}

\bibitem[\citeproctext]{ref-helleiner2025}
\CSLLeftMargin{4. }%
\CSLRightInline{E. Helleiner, M. DiLeo, J. van 't Klooster,
\href{https://doi.org/10.1111/rego.12629}{Financial {Technocrats} as
{Competitive Regime Creators}: {The Founding} and {Design} of the
{Network} for {Greening} the {Financial System}}. \emph{Regulation \&
Governance} \textbf{Online First} (2025).}

\bibitem[\citeproctext]{ref-garriga2025}
\CSLLeftMargin{5. }%
\CSLRightInline{A. C. Garriga,
\href{https://doi.org/10.1093/isq/sqaf024}{Revisiting {Central Bank
Independence} in the {World}: {An Extended Dataset}}. \emph{Int Stud Q}
\textbf{69}, sqaf024 (2025).}

\bibitem[\citeproctext]{ref-romelli2022}
\CSLLeftMargin{6. }%
\CSLRightInline{D. Romelli,
\href{https://doi.org/10.1093/epolic/eiac011}{The {Political Economy} of
{Reforms} in {Central Bank Design}: {Evidence} from a {New Dataset}}.
\emph{Econ Policy} \textbf{37}, 641--688 (2022).}

\bibitem[\citeproctext]{ref-chesler2023}
\CSLLeftMargin{7. }%
\CSLRightInline{A. Chesler, D. Javeline, K. Peh, S. Scogin,
\href{https://doi.org/10.1162/glep_a_00710}{Is {Democracy} the {Answer}
to {Intractable Climate Change}?} \emph{Global Environmental Politics}
\textbf{23}, 201--216 (2023).}

\bibitem[\citeproctext]{ref-sepulveda2025}
\CSLLeftMargin{8. }%
\CSLRightInline{M. V. Sepulveda, J. Piburn, The World Bank, Wbstats:
{Programmatic Access} to the {World Bank API}, (2025);
\url{https://cran.r-project.org/package=wbstats}.}

\bibitem[\citeproctext]{ref-mcgowan2025}
\CSLLeftMargin{9. }%
\CSLRightInline{C. McGowan, \emph{Describedata: {Miscellaneous
Descriptive Functions}} (2025;
\url{https://CRAN.R-project.org/package=describedata}).}

\bibitem[\citeproctext]{ref-ostrom1990}
\CSLLeftMargin{10. }%
\CSLRightInline{E. Ostrom, \emph{Governing the {Commons}: {The
Evolution} of {Institutions} for {Collective Action}} (Cambridge
University Press, New York, 1990).}

\bibitem[\citeproctext]{ref-maerz2026}
\CSLLeftMargin{11. }%
\CSLRightInline{S. Maerz, A. Edgell, S. Hellmeier, N. Ilchenko, L. Fox,
Vdemdata - an {R} package to load, explore and work with the most recent
{V-Dem} ({Varieties} of {Democracy}) and {V-Party} datasets (2026).
\url{https://www.v-dem.net/en/}.}

\bibitem[\citeproctext]{ref-coppedge2026vdem}
\CSLLeftMargin{12. }%
\CSLRightInline{M. Coppedge, J. Gerring, C. H. Knutsen, S. I. Lindberg,
J. Teorell, D. Altman, F. Angiolillo, M. Bernhard, A. Cornell, M. S.
Fish, L. Fox, L. Gastaldi, H. Gjerløw, A. Glynn, A. Good God, A. Hicken,
K. Kinzelbach, J. Krusell, K. L. Marquardt, K. McMann, V. Mechkova, J.
Medzihorsky, A. Neundorf, P. Paxton, D. Pemstein, J. Pernes, J. von
Römer, B. Seim, R. Sigman, S.-E. Skaaning, J. Staton, A. Sundström, M.
Tannenberg, E. Tzelgov, Y. Wang, T. Wig, S. Wilson, D. Ziblatt, V-dem
{[}country-year/country-date{]} dataset V16, Varieties of Democracy
(V-Dem) Project (2026); \url{https://doi.org/10.23696/vdemds26}.}

\bibitem[\citeproctext]{ref-povitkina2018}
\CSLLeftMargin{13. }%
\CSLRightInline{M. Povitkina,
\href{https://doi.org/10.1080/09644016.2018.1444723}{The {Limits} of
{Democracy} in {Tackling Climate Change}}. \emph{Environmental Politics}
\textbf{27}, 411--432 (2018).}

\bibitem[\citeproctext]{ref-ostrom2005}
\CSLLeftMargin{14. }%
\CSLRightInline{E. Ostrom, \emph{Understanding {Institutional
Diversity}} (Princeton University Press, Princeton, 2005).}

\bibitem[\citeproctext]{ref-vaccaro2023}
\CSLLeftMargin{15. }%
\CSLRightInline{A. Vaccaro,
\href{https://doi.org/10.1007/s11135-022-01466-x}{Measures of {State
Capacity}: {So Similar}, yet so {Different}}. \emph{Qual Quant}
\textbf{57}, 2281--2302 (2023).}

\bibitem[\citeproctext]{ref-pepinsky2018}
\CSLLeftMargin{16. }%
\CSLRightInline{T. B. Pepinsky,
\href{https://doi.org/10.1017/pan.2018.18}{A {Note} on {Listwise
Deletion} versus {Multiple Imputation}}. \emph{Political Analysis}
\textbf{26}, 480--488 (2018).}

\end{CSLReferences}




\end{document}
